\section*{Chapitre \ref{ch:b2:frontier-orbitals} --- Orbitales frontières et réactivité}

\begin{solution}{exo:b2:frontier-orbitals:1}
Éthène~: \omterm{def:b2:frontier-orbitals:frontier}{HOMO} $\alpha + \beta$, \omterm{def:b2:frontier-orbitals:frontier}{LUMO} $\alpha - \beta$. Anion allyle (quatre électrons)~: \omterm{def:b2:frontier-orbitals:frontier}{HOMO} $\alpha$ (le
niveau non liant), \omterm{def:b2:frontier-orbitals:frontier}{LUMO} $\alpha - 1.414\beta$. Butadiène~: \omterm{def:b2:frontier-orbitals:frontier}{HOMO} $\alpha + 0.618\beta$, \omterm{def:b2:frontier-orbitals:frontier}{LUMO} $\alpha -
0.618\beta$.
\end{solution}

\begin{solution}{exo:b2:frontier-orbitals:2}
Durs~: \ce{OH-}, \ce{H+}, \ce{F-}. Mous~: \ce{I-}, \ce{RS-}, le carbone de l'iodométhane.
\end{solution}

\begin{solution}{exo:b2:frontier-orbitals:3}
Le butadiène (conformère s-cis accessible) et le cyclopentadiène (bloqué en s-cis) réagissent~;
le penta-1,4-diène n'est pas conjugué et ne réagit pas. Le \omterm{def:b2:frontier-orbitals:diels-alder}{diène} (2\textit{Z},4\textit{Z}) ne peut
atteindre la forme s-cis qu'en poussant l'un contre l'autre ses deux groupes méthyle internes~: il
réagit très mal.
\end{solution}

\begin{solution}{exo:b2:frontier-orbitals:4}
Cyclohex-3-ène-1-carbonitrile~: un cycle à six chaînons, la double liaison entre les anciens C2
et C3 du \omterm{def:b2:frontier-orbitals:diels-alder}{diène}, le groupe \ce{C#N} porté par l'un des deux carbones du cycle issus du
\omterm{def:b2:frontier-orbitals:diels-alder}{diénophile}.
\end{solution}

\begin{solution}{exo:b2:frontier-orbitals:5}
$\Delta\varepsilon = \qty{8}{eV}$~: stabilisation perturbative $2\beta^2/\Delta\varepsilon = \qty{0.25}{eV}$. Exacte~: niveau inférieur
$-6 - \sqrt{16 + 1} = \qty{-10.123}{eV}$, gain $2 \times 0.123 = \qty{0.246}{eV}$.
\end{solution}

\begin{solution}{exo:b2:frontier-orbitals:6}
L'iodométhane est un \omterm{def:b2:frontier-orbitals:hard-soft}{électrophile mou}~: \omterm{def:b2:frontier-orbitals:control}{contrôle orbitalaire}, attaque par l'atome de plus grand
coefficient dans la \omterm{def:b2:frontier-orbitals:frontier}{HOMO}, le carbone. Un carbocation libre est un \omterm{def:b2:frontier-orbitals:hard-soft}{électrophile dur}~: \omterm{def:b2:frontier-orbitals:control}{contrôle de charge},
attaque en partie par l'azote, qui porte davantage de charge négative.
\end{solution}

\begin{solution}{exo:b2:frontier-orbitals:7}
Le C4 du \omterm{def:b2:frontier-orbitals:diels-alder}{diène} (plus grand coefficient dans la \omterm{def:b2:frontier-orbitals:frontier}{HOMO}) se lie au carbone $\beta$ du propénal (plus grand
coefficient dans la \omterm{def:b2:frontier-orbitals:frontier}{LUMO})~: 2-méthoxycyclohex-3-ène-1-carbaldéhyde, les deux substituants sur
des carbones voisins (produit «~ortho~»).
\end{solution}

\begin{solution}{exo:b2:frontier-orbitals:8}
Un squelette bicyclo[2.2.1]heptène (norbornène) portant le groupe ester sur un carbone de
l'ancien \omterm{def:b2:frontier-orbitals:diels-alder}{diénophile}, orienté vers la nouvelle liaison \ce{C=C}, sous le cycle à six chaînons
(endo), du côté opposé au pont \ce{CH2}.
\end{solution}

\begin{solution}{exo:b2:frontier-orbitals:9}
Les deux groupes carbonyle abaissent la \omterm{def:b2:frontier-orbitals:frontier}{LUMO} du \omterm{def:b2:frontier-orbitals:diels-alder}{diénophile}~: son écart avec la \omterm{def:b2:frontier-orbitals:frontier}{HOMO} du
cyclopentadiène est bien plus petit que celui de l'éthène, la stabilisation $\beta^2/\Delta\varepsilon$ de
l'état de transition est plus grande, et la barrière plus basse.
\end{solution}

\begin{solution}{exo:b2:frontier-orbitals:10}
La réaction est stéréospécifique~: le maléate donne l'adduit dont les deux groupes ester sont cis
(endo,endo comme produit majoritaire~; il est achiral, c'est un composé méso)~; le fumarate donne
l'adduit trans, un ester endo et l'autre exo, formé sous forme d'un couple d'énantiomères (racémique).
\end{solution}

\begin{solution}{exo:b2:frontier-orbitals:11}
Deux liaisons $\pi$ deviennent deux liaisons $\sigma$ plus fortes~: $\Delta_r H^\circ < 0$. Deux molécules n'en font plus qu'une~:
$\Delta_r S^\circ < 0$. $\Delta_r G^\circ = \Delta_r H^\circ - T\Delta_r S^\circ$ croît avec $T$ et devient
positive au-dessus de $T_i = \Delta_r H^\circ/\Delta_r S^\circ$~: à haute température, l'adduit se scinde de nouveau
(rétro-Diels--Alder).
\end{solution}

\begin{solution}{exo:b2:frontier-orbitals:12}
Deux orbitales de doublets non liants pleines côte à côte forment une interaction à quatre électrons, qui
est répulsive. La molécule tourne autour de la liaison \ce{N-N} jusqu'à ce que les doublets soient à peu près
perpendiculaires (une conformation décalée gauche), où leur recouvrement, et donc la répulsion, est
faible.
\end{solution}

\begin{solution}{pb:b2:frontier-orbitals:1}
\textbf{1.} Un cycle à cinq chaînons portant deux doubles liaisons conjuguées et un \ce{CH2}~: le cycle
bloque le motif \omterm{def:b2:frontier-orbitals:diels-alder}{diène} dans la conformation s-cis.
\textbf{2.} Le \omterm{def:b2:frontier-orbitals:diels-alder}{diénophile}, par l'une de ses doubles liaisons.
\textbf{3.} Un squelette norbornène accolé à un cycle cyclopentène~; les deux nouvelles liaisons $\sigma$ relient
les extrémités du \omterm{def:b2:frontier-orbitals:diels-alder}{diène} aux deux carbones de la double liaison du \omterm{def:b2:frontier-orbitals:diels-alder}{diénophile}.
\textbf{4.} L'\omterm{def:b2:frontier-orbitals:endo}{adduit endo} (\omterm{def:b2:frontier-orbitals:endo}{règle endo})~: le cycle restant du \omterm{def:b2:frontier-orbitals:diels-alder}{diénophile} se trouve sous
le \omterm{def:b2:frontier-orbitals:diels-alder}{diène} dans l'état de transition.
\textbf{5.} Oui~: elle est concertée, les deux liaisons se forment sur la même face de chaque partenaire, et la
géométrie du \omterm{def:b2:frontier-orbitals:diels-alder}{diénophile} est conservée.
\textbf{6.} Le cyclopentadiène est à la fois un bon \omterm{def:b2:frontier-orbitals:diels-alder}{diène} (bloqué en s-cis) et un \omterm{def:b2:frontier-orbitals:diels-alder}{diénophile}, et la
réaction est concertée, sans intermédiaire chargé à stabiliser.
\textbf{7.} Négatif~: deux liaisons $\pi$ sont remplacées par deux liaisons $\sigma$.
\textbf{8.} Négatif~: deux molécules n'en donnent qu'une.
\textbf{9.} $\Delta_r G^\circ = \Delta_r H^\circ - T\Delta_r S^\circ$, les deux termes étant négatifs~: négative à basse
$T$, nulle en $T_i = \Delta_r H^\circ/\Delta_r S^\circ$, positive au-delà.
\textbf{10.} À température ambiante, $T < T_i$~: la dimérisation est favorisée, bien que lente. Près
de la température d'ébullition du dimère, l'équilibre favorise le monomère, et il est atteint
rapidement.
\textbf{11.} Le monomère quitte le mélange chaud sous forme de vapeur~: retirer un produit maintient le
\omterm{def:b2:equilibrium-shifts:shift}{déplacement de l'équilibre} vers le monomère (Le Chatelier).
\textbf{12.} À température ambiante, il se dimérise de nouveau~; le froid ralentit cette réaction.
\textbf{13.} $0.62 - (-0.62) = 1.24|\beta|$.
\textbf{14.} $0.62 - (-0.35) = 0.97|\beta|$.
\textbf{15.} Le cyclopentadiène avec l'anhydride maléique~: le plus petit écart donne une plus grande
stabilisation de l'état de transition, donc une réaction plus rapide.
\textbf{16.} Ils sont conjugués avec la liaison \ce{C=C} et attracteurs d'électrons~: comme dans le
propénal, ils abaissent le niveau $\pi^*$.
\textbf{17.} Sa \omterm{def:b2:frontier-orbitals:frontier}{LUMO}, qui présente des lobes sur les carbones des carbonyles, face aux carbones centraux de
la \omterm{def:b2:frontier-orbitals:frontier}{HOMO} du \omterm{def:b2:frontier-orbitals:diels-alder}{diène}.
\textbf{18.} Squelette norbornène~; le cycle anhydride sous la nouvelle double liaison, du côté opposé au
pont \ce{CH2}.
\textbf{19.} Les deux carbones en tête de pont et les deux carbones qui portent l'anhydride~: quatre
stéréocentres.
\textbf{20.} Non~: un plan de symétrie passe par le pont \ce{CH2} et entre les deux moitiés
de la molécule~; c'est un composé méso.
\textbf{21.} Ils étaient cis sur le \omterm{def:b2:frontier-orbitals:diels-alder}{diénophile}, et l'addition concertée les maintient cis.
\textbf{22.} Le même squelette avec le cycle anhydride du côté du pont \ce{CH2}.
\textbf{23.} Non~: passer de endo à exo exigerait de rompre et de reformer des liaisons \ce{C-C}, par
la réaction inverse, qui demande un chauffage.
\textbf{24.} L'eau ouvre l'anhydride~: on obtient le diacide carboxylique cis, les deux groupes \ce{COOH} étant endo.
\textbf{25.} L'\omterm{def:b2:frontier-orbitals:endo}{adduit endo} possède $\boldsymbol{4}$ stéréocentres.
\end{solution}
